\section{Checking partial EEG epoch replay beyond complete recording identity}
\label{sec:v27-epoch-identity}
The V26 check identified one pair of identical full
preprocessed EEG recording arrays under different subject identifiers.
That check is insufficient to exclude \emph{partial} record overlap:
two different 40-trial arrays can share a small number of identical
epochs while having different complete-record SHA-256 digests.
We distinguish EEG \emph{subject identifiers}, numerical
\emph{recording content}, and human patient identity; the last
cannot be determined from the published arrays.

\paragraph{Content-based epoch-overlap graph.}
For subject identifier $i$, let
$E_i=\{h(x_{it}):t=1,\ldots,n_i\}$ be the set of
cryptographic fingerprints of the complete EEG trial arrays
after the single, frozen MOABB preprocessing pipeline. The
fingerprint hashes both channel/time dimensions and the contiguous
float64 signal bytes. Candidate equality is checked by exact
numerical array comparison, with contradictory class labels
treated as a dataset-integrity failure.
Define the undirected graph with an edge $(i,j)$ if
$E_i\cap E_j\ne\varnothing$ for $i\ne j$.
The graph ignores observed classifier predictions and AUC values.

For any partition that assigns each complete identifier's
trials to one fold, \emph{zero exact trial-content overlap between
training and test folds} is equivalent to requiring that
every connected component of this graph is contained in
one fold. If an edge crosses the partition, its identical
epoch appears on both sides; conversely, if an exact epoch
appears on both sides, that pair defines a crossing edge.
A component can be separated only by cutting such an edge.
This elementary graph observation is not a newly invented
generalisation theorem. It gives an actionable grouping
contract under exact equality, which the original
complete-record check does not provide.

\paragraph{Audit and boundary.}
We precommitted a retrospective scan over all 50 public
Liu2024 subject IDs and all 2,000 originally retained
40-trial observations, before executing the new trial-level
scan. A pinned MOABB 1.7.2 loader retains the previous
29-channel, 8--30 Hz, 160 Hz left/right motor-imagery
processing and 20/20 class balance. For each exact epoch
content appearing under multiple identifiers, we count
shared unique epoch contents, cross-identifier matching
trial instances and any transitive groups. Internal repeats
are reported separately. No threshold is tuned to target
classifier performance, and no unverified near-duplicate
criterion is introduced. Synthetic tests inject
partial overlap, transitive reuse, whole-record copies,
and conflicting event labels. The entire source EEG
trial list and matching rules are fixed prior to the
audit, with complete reproducible logs archived.

All grouping necessarily inspects globally released
EEG signal content, including held records. Consequently,
any resulting component-aware reanalysis is an
\emph{exposed-cohort integrity sensitivity study},
not a physically target-isolated, independently
prospective zero-calibration experiment.
Exact equality says nothing about near duplicates,
overlapping-but-nonidentical time windows, patient
identity or source-site dependencies. Even a negative
additional-overlap scan would only strengthen one
specific data-integrity check and cannot be promoted
to a new AUC gain or proof of population transfer.

Published EEG validation research documents
the importance of the subject/record split
\cite{gil2026subjectvalidation}. Our novelty claim
is restricted to the \emph{specific} released-record
finding and its reproducible trial-granularity
consequences under a frozen EEG decoding pipeline,
not to the general concept of preventing leakage.
